\section{Formal Methodology}
\label{sec:formal_methodology}

\begin{BLACKSHEET}

To verify the MCAS algorithm, we developed a formal specification in TLA+ suitable for the TLC model checker. The model abstracts the system into a state machine where the shared memory mem is represented as a record containing two fields: mem.val, storing the current integer value, and mem.owner, which holds either a NULL value or the ID of the thread currently holding an exclusive lock on the cell. To capture asynchronous execution, each thread t is modeled with its own control state (pc[t]), an announcement flag (announce), and a transaction status (status) indicating whether its operations are Undecided, Committed, or Aborted.

\textbf{Need to write the part with description of specification for other fields}
A. Writer Lifecycle: Acquisition and Atomic Commit
....
\textbf{Need to write the part with description of specification for the writer}
B. Reader Lifecycle: Announcement and Validation
....
\textbf{Need to write the part with description of specification for the reader}
C. Formal Verification and Correctness We verify the MCAS model as a state-transition system using the TLA+ framework. The safety of the algorithm is established through the following invariants:
Safety Invariant (linearisability...):
System Progress (lock-free...):
\textbf{Need to write the part with description of verification}
....

\end{BLACKSHEET}
